Many current agent architectures layer retrieval, memory, orchestration, and infrastructure as loosely coupled modules. A vector store handles semantic lookup, a graph may be added as an auxiliary index, and runtime resource constraints are usually treated as implementation details rather than part of the model. This separation is often practical, but it limits what the system can explain about its own behavior.

This paper explores a stronger systems hypothesis: semantics, topology, provenance, and resource pressure should be represented inside the same runtime substrate. In the proposed design, all persisted observations live in an append-only \truthgraph{}, while computation proceeds over a compressed and inspectable \viewgraph{}. Decision-making entities, called presences, emit and absorb probabilistic packets called daimoi. Routing is shaped by graph distance, semantic affinity, congestion, and local scarcity-derived prices.

The result is intended to be more than a graph-enhanced RAG stack. It is a runtime in which semantic organization affects path structure, path structure affects transport cost, transport cost affects behavior, and all of these effects are tied back to ledgered events. The core promise is therefore not only capability but auditability.

The draft makes four concrete claims. First, it argues for a \truthgraph{}/\viewgraph{} split in which immutable truth and operational memory remain connected but not collapsed. Second, it introduces presences and daimoi as explicit control-layer objects rather than treating all routing as implicit prompt orchestration. Third, it treats compute and IO pressure as first-class graph state via a reservoir economy instead of hiding them behind the runtime. Fourth, it proposes an evaluation style based on workloads, ablations, and audit trails rather than a generic leaderboard.

This draft focuses on the paper boundary that seems strongest: the graph-native runtime itself. The wider \texttt{fork\_tales} repository contains more ideas than one paper should carry, so the present draft intentionally prioritizes the substrate claim over the larger mythic or ecological framing.
